
% --- SLIDE 4: OPTUNA CODE ---
\begin{frame}[fragile]{Optuna: Define-by-Run}

    \begin{lstlisting}[language=Python, basicstyle=\ttfamily\scriptsize]
import optuna

def objective(trial):
    # the search space is *built by running Python*, so it can branch
    model_name = trial.suggest_categorical("model", ["rf", "gbm"])
    if model_name == "rf":
        clf = RandomForestClassifier(
            n_estimators = trial.suggest_int("n_estimators", 100, 1000, step=100),
            max_depth    = trial.suggest_int("max_depth", 2, 32, log=True),
            random_state = 42)
    else:
        clf = HistGradientBoostingClassifier(
            learning_rate = trial.suggest_float("lr", 1e-3, 3e-1, log=True),
            max_leaf_nodes= trial.suggest_int("leaves", 8, 256, log=True),
            random_state  = 42)
    return cross_val_score(clf, X, y, cv=5, scoring="f1_macro").mean()

study = optuna.create_study(direction="maximize")   # default sampler: TPE
study.optimize(objective, n_trials=100)
print(study.best_params, study.best_value)
    \end{lstlisting}

    \vspace{0.09em}
    \begin{itemize}
        \scriptsize
        \item \textbf{Define-by-run} is Optuna's headline design principle: the search space is constructed
              dynamically as the objective executes, so an \texttt{if} statement gives you a conditional
              space for free.
        \item \texttt{log=True} for anything spanning orders of magnitude (learning rate, regularisation
              strength). Sampling learning rate uniformly in $[10^{-5}, 10^{-1}]$ wastes 99\% of your budget
              above $10^{-2}$.
    \end{itemize}

\end{frame}
